% \begin{table}[t]
%     \centering
%     \caption[نتیجه‌ی اعمال الگوریتم شناسایی لایه‌ی بهینه در مجموعه داده‌ی \lr{Cifar100}.]{نتایج اعمال الگوریتم شناسایی لایه‌ی بهینه بر روی مجموعه‌داده‌ی \lr{CIFAR100} و مدل مبتنی بر مبدل بینایی. در این مدل، لایه‌ی ۰ متناظر با خروجی لایه‌ی \lr{Patch Embedding} و دوازده لایه‌ی بعدی متناظر با خروجی‌ لایه‌های دوازده \lr{Transformer Encoder} هستند. بر اساس مقادیر امتیاز جرم تقریبی، لایه‌ی چهارم به‌عنوان لایه‌ی بهینه انتخاب می‌شود.}    \label{table:alg2_results}
%     \begin{latin}
%     \begin{tabular}{c c}
%         \toprule
%         Layer Index & APMS \\
%         \midrule
%         0  & 0.0117 \\
%         1  & 0.0148 \\
%         2  & 0.0190 \\
%         3  & 0.0406 \\
%         4  & 0.0508 \\
%         5  & 0.0496 \\
%         6  & 0.0440 \\
%         7  & 0.0368 \\
%         8  & 0.0246 \\
%         9  & 0.0196 \\
%         10 & 0.0039 \\
%         11 & 0.0024 \\
%         12 & 0.0272 \\
%         \bottomrule
%     \end{tabular}
%     \end{latin}
% \end{table}
\begin{table}[t]
    \centering
    \caption[نتیجه‌ی اعمال الگوریتم شناسایی لایه‌ی بهینه در مجموعه داده‌ی \lr{Cifar100}.]{نتایج اعمال الگوریتم شناسایی لایه‌ی بهینه بر روی مجموعه‌داده‌ی \lr{CIFAR100} و مدل مبتنی بر مبدل بینایی. در این مدل، لایه‌ی ۰ متناظر با خروجی لایه‌ی \lr{Patch Embedding} و دوازده لایه‌ی بعدی متناظر با خروجی‌ لایه‌های دوازده \lr{Transformer Encoder} هستند. بر اساس مقادیر امتیاز جرم تقریبی، لایه‌ی چهارم به‌عنوان لایه‌ی بهینه انتخاب می‌شود.}    
    \label{table:alg2_results}
    \begin{latin}
    \resizebox{\textwidth}{!}{
    \begin{tabular}{c c c c c c c c c c c c c c}
        \toprule
        Layer Index & 0 & 1 & 2 & 3 & 4 & 5 & 6 & 7 & 8 & 9 & 10 & 11 & 12 \\
        \midrule
        APMS        & 0.0117 & 0.0148 & 0.0190 & 0.0406 & 0.0508 & 0.0496 & 0.0440 & 0.0368 & 0.0246 & 0.0196 & 0.0039 & 0.0024 & 0.0272 \\
        \bottomrule
    \end{tabular}
    }
    \end{latin}
\end{table}
